\section{Anexos}
\label{sec:anexos}

En esta seccion se incluye el codigo fuente completo del Modelo matematico y la salida completa de la validacion manual de la Seccion~\ref{sec:validacion}. El codigo corresponde a la carpeta \texttt{src/BezierComparador/Modelo} del repositorio.

\subsection{Codigo Fuente del Modelo Matematico}
\label{subsec:anx_codigo}

El Listado~\ref{lst:anx_bernstein} implementa la Ecuacion~\ref{eq:bernstein} con coeficientes binomiales precalculados. El Listado~\ref{lst:anx_casteljau} implementa la Ecuacion~\ref{eq:casteljau} y el metodo \texttt{ObtenerNivelesIntermedios}, que usa la Vista para dibujar las interpolaciones. El Listado~\ref{lst:anx_comparador} calcula las diferencias numericas entre ambos metodos y el Listado~\ref{lst:anx_curva} administra los puntos de control.

\begin{lstlisting}[language={[Sharp]C}, caption={Clase clsEvaluadorBernstein (evaluacion con polinomios de Bernstein).}, label={lst:anx_bernstein}]
using System;
using System.Collections.Generic;

namespace BezierComparador
{
    // Evalua la curva de Bezier en forma directa con polinomios de Bernstein
    // Complejidad: O(n) por muestra, O(N * n) para N muestras
    public class clsEvaluadorBernstein : IEvaluadorCurva
    {
        public string CadNombre
        {
            get { return "Bernstein"; }
        }

        // Formula: B(t) = suma para i=0..n de C(n,i) * (1-t)^(n-i) * t^i * P_i
        // Parametro: dblT = k / (entMuestras - 1), con k de 0 a entMuestras - 1
        // Los coeficientes C(n,i) se calculan una sola vez, fuera del ciclo de muestras
        public clsResultadoEvaluacion Evaluar(List<clsPuntoControl> lstPuntos, int entMuestras)
        {
            clsResultadoEvaluacion objResultado = new clsResultadoEvaluacion(CadNombre);

            // Valida minimo 2 puntos de control y 2 muestras
            if (lstPuntos == null || lstPuntos.Count < clsCurvaBezier.entMinimoPuntos || entMuestras < 2)
            {
                return objResultado;
            }

            int entN = lstPuntos.Count - 1;
            double[] arrCoeficientes = CalcularCoeficientesBinomiales(entN);
            clsPuntoControl[] arrCurva = new clsPuntoControl[entMuestras];

            // Recorre cada muestra del parametro t entre 0 y 1
            for (int entK = 0; entK < entMuestras; entK++)
            {
                double dblT = (double)entK / (double)(entMuestras - 1);
                double dblSumaX = 0.0;
                double dblSumaY = 0.0;

                // Suma ponderada de los puntos de control: n + 1 terminos, O(n)
                for (int entI = 0; entI <= entN; entI++)
                {
                    double dblBase = CalcularBase(arrCoeficientes[entI], entN, entI, dblT);
                    dblSumaX += dblBase * lstPuntos[entI].dblX;
                    dblSumaY += dblBase * lstPuntos[entI].dblY;
                }

                arrCurva[entK] = new clsPuntoControl(dblSumaX, dblSumaY);
            }

            objResultado.arrPuntosCurva = arrCurva;
            return objResultado;
        }

        // Formula: C(n,0) = 1 y C(n,i) = C(n,i-1) * (n - i + 1) / i
        // Iterativo sin factorial para evitar desbordamiento, O(n)
        public double[] CalcularCoeficientesBinomiales(int entN)
        {
            double[] arrCoef = new double[entN + 1];
            arrCoef[0] = 1.0;

            for (int entI = 1; entI <= entN; entI++)
            {
                arrCoef[entI] = arrCoef[entI - 1] * (double)(entN - entI + 1) / (double)entI;
            }

            return arrCoef;
        }

        // Formula: B_{i,n}(t) = C(n,i) * (1-t)^(n-i) * t^i
        public double CalcularBase(double dblCoeficiente, int entN, int entI, double dblT)
        {
            double dblPotenciaUnoMenosT = ElevarPotencia(1.0 - dblT, entN - entI);
            double dblPotenciaT = ElevarPotencia(dblT, entI);
            return dblCoeficiente * dblPotenciaUnoMenosT * dblPotenciaT;
        }

        // Calcula base^exponente con Math.Pow, tomando 0^0 = 1
        // Asi en t = 0 y t = 1 la curva pasa exacto por P_0 y P_n
        private double ElevarPotencia(double dblBase, int entExponente)
        {
            if (entExponente == 0)
            {
                return 1.0;
            }
            return Math.Pow(dblBase, entExponente);
        }
    }
}
\end{lstlisting}

\begin{lstlisting}[language={[Sharp]C}, caption={Clase clsEvaluadorDeCasteljau (algoritmo de De Casteljau).}, label={lst:anx_casteljau}]
using System.Collections.Generic;

namespace BezierComparador
{
    // Evalua la curva de Bezier con interpolaciones lineales sucesivas (De Casteljau)
    // Complejidad: O(n^2) por muestra, O(N * n^2) para N muestras
    public class clsEvaluadorDeCasteljau : IEvaluadorCurva
    {
        public string CadNombre
        {
            get { return "De Casteljau"; }
        }

        // Formula: P_i^(r) = (1-t) * P_i^(r-1) + t * P_(i+1)^(r-1)
        // r de 1..n, i de 0..n-r. El punto de la curva es P_0^(n)
        // Parametro: dblT = k / (entMuestras - 1), con k de 0 a entMuestras - 1
        public clsResultadoEvaluacion Evaluar(List<clsPuntoControl> lstPuntos, int entMuestras)
        {
            clsResultadoEvaluacion objResultado = new clsResultadoEvaluacion(CadNombre);

            // Valida minimo 2 puntos de control y 2 muestras
            if (lstPuntos == null || lstPuntos.Count < clsCurvaBezier.entMinimoPuntos || entMuestras < 2)
            {
                return objResultado;
            }

            int entN = lstPuntos.Count - 1;
            clsPuntoControl[] arrCurva = new clsPuntoControl[entMuestras];

            // Arreglos de trabajo para la copia de X y Y (se reservan una sola vez)
            double[] arrX = new double[entN + 1];
            double[] arrY = new double[entN + 1];

            // Recorre cada muestra del parametro t entre 0 y 1
            for (int entK = 0; entK < entMuestras; entK++)
            {
                double dblT = (double)entK / (double)(entMuestras - 1);
                double dblUnoMenosT = 1.0 - dblT;

                // Copia los puntos de control (nivel r = 0) para no alterar la lista original
                for (int entI = 0; entI <= entN; entI++)
                {
                    arrX[entI] = lstPuntos[entI].dblX;
                    arrY[entI] = lstPuntos[entI].dblY;
                }

                // Reduce por niveles: cada nivel tiene un punto menos, n(n+1)/2 operaciones
                for (int entR = 1; entR <= entN; entR++)
                {
                    for (int entI = 0; entI <= entN - entR; entI++)
                    {
                        arrX[entI] = dblUnoMenosT * arrX[entI] + dblT * arrX[entI + 1];
                        arrY[entI] = dblUnoMenosT * arrY[entI] + dblT * arrY[entI + 1];
                    }
                }

                // Al final del ultimo nivel queda P_0^(n), el punto de la curva
                arrCurva[entK] = new clsPuntoControl(arrX[0], arrY[0]);
            }

            objResultado.arrPuntosCurva = arrCurva;
            return objResultado;
        }

        // Formula: P_i^(r) = (1-t) * P_i^(r-1) + t * P_(i+1)^(r-1) para un t fijo
        // Devuelve todos los niveles: nivel 0 = puntos de control, nivel n = punto de la curva
        // Se usa para dibujar las interpolaciones (no para medir tiempos)
        public List<List<clsPuntoControl>> ObtenerNivelesIntermedios(List<clsPuntoControl> lstPuntos, double dblT)
        {
            List<List<clsPuntoControl>> lstNiveles = new List<List<clsPuntoControl>>();

            if (lstPuntos == null || lstPuntos.Count < clsCurvaBezier.entMinimoPuntos)
            {
                return lstNiveles;
            }

            // Mantiene t dentro del intervalo [0, 1]
            if (dblT < 0.0)
            {
                dblT = 0.0;
            }
            if (dblT > 1.0)
            {
                dblT = 1.0;
            }
            double dblUnoMenosT = 1.0 - dblT;

            // Nivel 0: copia de los puntos de control
            List<clsPuntoControl> lstNivelCero = new List<clsPuntoControl>();
            for (int entI = 0; entI < lstPuntos.Count; entI++)
            {
                lstNivelCero.Add(lstPuntos[entI].Clonar());
            }
            lstNiveles.Add(lstNivelCero);

            // Niveles 1..n: interpola cada par de puntos consecutivos del nivel anterior
            int entN = lstPuntos.Count - 1;
            for (int entR = 1; entR <= entN; entR++)
            {
                List<clsPuntoControl> lstAnterior = lstNiveles[entR - 1];
                List<clsPuntoControl> lstNuevo = new List<clsPuntoControl>();

                for (int entI = 0; entI <= entN - entR; entI++)
                {
                    double dblX = dblUnoMenosT * lstAnterior[entI].dblX + dblT * lstAnterior[entI + 1].dblX;
                    double dblY = dblUnoMenosT * lstAnterior[entI].dblY + dblT * lstAnterior[entI + 1].dblY;
                    lstNuevo.Add(new clsPuntoControl(dblX, dblY));
                }

                lstNiveles.Add(lstNuevo);
            }

            return lstNiveles;
        }
    }
}
\end{lstlisting}

\begin{lstlisting}[language={[Sharp]C}, caption={Clase clsComparador (diferencias numericas entre ambos metodos).}, label={lst:anx_comparador}]
using System;

namespace BezierComparador
{
    // Calcula la diferencia numerica punto a punto entre dos curvas evaluadas
    public class clsComparador
    {
        // Formula por muestra: d_k = raiz((x1 - x2)^2 + (y1 - y2)^2)
        // Maxima: max(d_k). Promedio: (1/M) * suma(d_k)
        // Error cuadratico medio: ECM = (1/M) * suma(d_k^2)
        // M = cantidad de muestras comparadas. Complejidad O(M)
        public clsResultadoComparacion Comparar(clsPuntoControl[] arrCurvaA, clsPuntoControl[] arrCurvaB)
        {
            clsResultadoComparacion objResultado = new clsResultadoComparacion();

            if (arrCurvaA == null || arrCurvaB == null)
            {
                return objResultado;
            }

            // Si los tamanos difieren solo se comparan las muestras comunes
            int entCantidad = Math.Min(arrCurvaA.Length, arrCurvaB.Length);
            if (entCantidad == 0)
            {
                return objResultado;
            }

            double dblMaximo = 0.0;
            double dblSumaDistancias = 0.0;
            double dblSumaCuadrados = 0.0;

            // Recorre cada muestra k y acumula la distancia entre ambos puntos
            for (int entK = 0; entK < entCantidad; entK++)
            {
                double dblDifX = arrCurvaA[entK].dblX - arrCurvaB[entK].dblX;
                double dblDifY = arrCurvaA[entK].dblY - arrCurvaB[entK].dblY;
                double dblDistanciaCuadrada = dblDifX * dblDifX + dblDifY * dblDifY;
                double dblDistancia = Math.Sqrt(dblDistanciaCuadrada);

                if (dblDistancia > dblMaximo)
                {
                    dblMaximo = dblDistancia;
                }

                dblSumaDistancias += dblDistancia;
                dblSumaCuadrados += dblDistanciaCuadrada;
            }

            objResultado.dblDiferenciaMaxima = dblMaximo;
            objResultado.dblDiferenciaPromedio = dblSumaDistancias / entCantidad;
            objResultado.dblErrorCuadraticoMedio = dblSumaCuadrados / entCantidad;
            objResultado.entMuestrasComparadas = entCantidad;
            return objResultado;
        }

        // Compara directamente dos resultados de evaluacion
        public clsResultadoComparacion Comparar(clsResultadoEvaluacion objResultadoA, clsResultadoEvaluacion objResultadoB)
        {
            if (objResultadoA == null || objResultadoB == null)
            {
                return new clsResultadoComparacion();
            }
            return Comparar(objResultadoA.arrPuntosCurva, objResultadoB.arrPuntosCurva);
        }
    }
}
\end{lstlisting}

\begin{lstlisting}[language={[Sharp]C}, caption={Clase clsCurvaBezier (gestion de puntos de control y validaciones).}, label={lst:anx_curva}]
using System.Collections.Generic;

namespace BezierComparador
{
    // Almacena y gestiona los puntos de control de una curva de Bezier
    public class clsCurvaBezier
    {
        // Minimo de puntos para poder evaluar (curva de grado 1)
        public const int entMinimoPuntos = 2;

        public List<clsPuntoControl> lstPuntosControl;

        public List<clsPuntoControl> LstPuntosControl
        {
            get { return lstPuntosControl; }
            set { lstPuntosControl = value; }
        }

        public clsCurvaBezier()
        {
            lstPuntosControl = new List<clsPuntoControl>();
        }

        // Formula: grado n = cantidad de puntos - 1
        public int ObtenerGrado()
        {
            return lstPuntosControl.Count - 1;
        }

        // Devuelve cuantos puntos de control tiene la curva
        public int ObtenerCantidadPuntos()
        {
            return lstPuntosControl.Count;
        }

        // Verifica que el indice exista dentro de la lista
        public bool EsIndiceValido(int entIndice)
        {
            return entIndice >= 0 && entIndice < lstPuntosControl.Count;
        }

        // Verifica que haya al menos 2 puntos para evaluar
        public bool PuedeEvaluar()
        {
            return lstPuntosControl.Count >= entMinimoPuntos;
        }

        // Mensaje para el usuario si la curva no se puede evaluar
        public string ObtenerMensajeValidacion()
        {
            if (PuedeEvaluar())
            {
                return "";
            }
            return "Agregue al menos " + entMinimoPuntos + " puntos";
        }

        // Agrega un punto al final; devuelve false si es nulo
        public bool AgregarPunto(clsPuntoControl ptoNuevo)
        {
            if (ptoNuevo == null)
            {
                return false;
            }
            lstPuntosControl.Add(ptoNuevo);
            return true;
        }

        // Agrega un punto con coordenadas X y Y
        public bool AgregarPunto(double dblX, double dblY)
        {
            return AgregarPunto(new clsPuntoControl(dblX, dblY));
        }

        // Devuelve el punto del indice o null si el indice no existe
        public clsPuntoControl? ObtenerPunto(int entIndice)
        {
            if (!EsIndiceValido(entIndice))
            {
                return null;
            }
            return lstPuntosControl[entIndice];
        }

        // Elimina el punto del indice; devuelve false si el indice no existe
        public bool EliminarPunto(int entIndice)
        {
            if (!EsIndiceValido(entIndice))
            {
                return false;
            }
            lstPuntosControl.RemoveAt(entIndice);
            return true;
        }

        // Cambia las coordenadas del punto; devuelve false si el indice no existe
        public bool ModificarPunto(int entIndice, double dblX, double dblY)
        {
            if (!EsIndiceValido(entIndice))
            {
                return false;
            }
            lstPuntosControl[entIndice].dblX = dblX;
            lstPuntosControl[entIndice].dblY = dblY;
            return true;
        }

        // Elimina todos los puntos de control
        public void Limpiar()
        {
            lstPuntosControl.Clear();
        }
    }
}
\end{lstlisting}

% Andres y Joseph: agregar aqui como \subsection el codigo completo de sus clases
% (clsMedidor, clsControladorPruebas, clsControladorCurva, ctlLienzo, etc.)
% con el mismo formato lstlisting y etiquetas lst:anx_<nombre>.

\subsection{Salida Completa de la Validacion Manual}
\label{subsec:anx_salida}

El Listado~\ref{lst:anx_salida} muestra la salida del proyecto de consola \texttt{BezierComparador.PruebaManual}, guardada en el archivo \texttt{validacion\_manual.txt}. Para regenerarla se ejecuta \texttt{dotnet run --project src/BezierComparador.PruebaManual} desde la raiz del repositorio.

\begin{lstlisting}[caption={Salida de la validacion manual de la curva cubica.}, label={lst:anx_salida}, basicstyle=\ttfamily\scriptsize]
VALIDACION MATEMATICA MANUAL - CURVA CUBICA (n = 3)
Puntos de control:
  P0 = (0.000000, 0.000000)
  P1 = (1.000000, 3.000000)
  P2 = (4.000000, 3.000000)
  P3 = (5.000000, 0.000000)

Coeficientes binomiales C(3,i): 1 3 3 1 

==================================================
t = 0.00   (1 - t = 1.00)
==================================================
[Bernstein]  B(t) = suma C(n,i) * (1-t)^(n-i) * t^i * P_i
  i=0: B0,3 = 1 * 1.00^3 * 0.00^0 = 1.000000  ->  B * P0 = (0.000000, 0.000000)
  i=1: B1,3 = 3 * 1.00^2 * 0.00^1 = 0.000000  ->  B * P1 = (0.000000, 0.000000)
  i=2: B2,3 = 3 * 1.00^1 * 0.00^2 = 0.000000  ->  B * P2 = (0.000000, 0.000000)
  i=3: B3,3 = 1 * 1.00^0 * 0.00^3 = 0.000000  ->  B * P3 = (0.000000, 0.000000)
  Suma: B(0.00) = (0.000000, 0.000000)
[De Casteljau]  P_i^(r) = (1-t) * P_i^(r-1) + t * P_(i+1)^(r-1)
  r=0: P0^(0)=(0.000000, 0.000000)  P1^(0)=(1.000000, 3.000000)  P2^(0)=(4.000000, 3.000000)  P3^(0)=(5.000000, 0.000000)  
  r=1: P0^(1)=(0.000000, 0.000000)  P1^(1)=(1.000000, 3.000000)  P2^(1)=(4.000000, 3.000000)  
  r=2: P0^(2)=(0.000000, 0.000000)  P1^(2)=(1.000000, 3.000000)  
  r=3: P0^(3)=(0.000000, 0.000000)  
  Resultado: C(0.00) = (0.000000, 0.000000)
Diferencia |Bernstein - De Casteljau| = 0.000E+000

==================================================
t = 0.25   (1 - t = 0.75)
==================================================
[Bernstein]  B(t) = suma C(n,i) * (1-t)^(n-i) * t^i * P_i
  i=0: B0,3 = 1 * 0.75^3 * 0.25^0 = 0.421875  ->  B * P0 = (0.000000, 0.000000)
  i=1: B1,3 = 3 * 0.75^2 * 0.25^1 = 0.421875  ->  B * P1 = (0.421875, 1.265625)
  i=2: B2,3 = 3 * 0.75^1 * 0.25^2 = 0.140625  ->  B * P2 = (0.562500, 0.421875)
  i=3: B3,3 = 1 * 0.75^0 * 0.25^3 = 0.015625  ->  B * P3 = (0.078125, 0.000000)
  Suma: B(0.25) = (1.062500, 1.687500)
[De Casteljau]  P_i^(r) = (1-t) * P_i^(r-1) + t * P_(i+1)^(r-1)
  r=0: P0^(0)=(0.000000, 0.000000)  P1^(0)=(1.000000, 3.000000)  P2^(0)=(4.000000, 3.000000)  P3^(0)=(5.000000, 0.000000)  
  r=1: P0^(1)=(0.250000, 0.750000)  P1^(1)=(1.750000, 3.000000)  P2^(1)=(4.250000, 2.250000)  
  r=2: P0^(2)=(0.625000, 1.312500)  P1^(2)=(2.375000, 2.812500)  
  r=3: P0^(3)=(1.062500, 1.687500)  
  Resultado: C(0.25) = (1.062500, 1.687500)
Diferencia |Bernstein - De Casteljau| = 0.000E+000

==================================================
t = 0.50   (1 - t = 0.50)
==================================================
[Bernstein]  B(t) = suma C(n,i) * (1-t)^(n-i) * t^i * P_i
  i=0: B0,3 = 1 * 0.50^3 * 0.50^0 = 0.125000  ->  B * P0 = (0.000000, 0.000000)
  i=1: B1,3 = 3 * 0.50^2 * 0.50^1 = 0.375000  ->  B * P1 = (0.375000, 1.125000)
  i=2: B2,3 = 3 * 0.50^1 * 0.50^2 = 0.375000  ->  B * P2 = (1.500000, 1.125000)
  i=3: B3,3 = 1 * 0.50^0 * 0.50^3 = 0.125000  ->  B * P3 = (0.625000, 0.000000)
  Suma: B(0.50) = (2.500000, 2.250000)
[De Casteljau]  P_i^(r) = (1-t) * P_i^(r-1) + t * P_(i+1)^(r-1)
  r=0: P0^(0)=(0.000000, 0.000000)  P1^(0)=(1.000000, 3.000000)  P2^(0)=(4.000000, 3.000000)  P3^(0)=(5.000000, 0.000000)  
  r=1: P0^(1)=(0.500000, 1.500000)  P1^(1)=(2.500000, 3.000000)  P2^(1)=(4.500000, 1.500000)  
  r=2: P0^(2)=(1.500000, 2.250000)  P1^(2)=(3.500000, 2.250000)  
  r=3: P0^(3)=(2.500000, 2.250000)  
  Resultado: C(0.50) = (2.500000, 2.250000)
Diferencia |Bernstein - De Casteljau| = 0.000E+000

==================================================
t = 0.75   (1 - t = 0.25)
==================================================
[Bernstein]  B(t) = suma C(n,i) * (1-t)^(n-i) * t^i * P_i
  i=0: B0,3 = 1 * 0.25^3 * 0.75^0 = 0.015625  ->  B * P0 = (0.000000, 0.000000)
  i=1: B1,3 = 3 * 0.25^2 * 0.75^1 = 0.140625  ->  B * P1 = (0.140625, 0.421875)
  i=2: B2,3 = 3 * 0.25^1 * 0.75^2 = 0.421875  ->  B * P2 = (1.687500, 1.265625)
  i=3: B3,3 = 1 * 0.25^0 * 0.75^3 = 0.421875  ->  B * P3 = (2.109375, 0.000000)
  Suma: B(0.75) = (3.937500, 1.687500)
[De Casteljau]  P_i^(r) = (1-t) * P_i^(r-1) + t * P_(i+1)^(r-1)
  r=0: P0^(0)=(0.000000, 0.000000)  P1^(0)=(1.000000, 3.000000)  P2^(0)=(4.000000, 3.000000)  P3^(0)=(5.000000, 0.000000)  
  r=1: P0^(1)=(0.750000, 2.250000)  P1^(1)=(3.250000, 3.000000)  P2^(1)=(4.750000, 0.750000)  
  r=2: P0^(2)=(2.625000, 2.812500)  P1^(2)=(4.375000, 1.312500)  
  r=3: P0^(3)=(3.937500, 1.687500)  
  Resultado: C(0.75) = (3.937500, 1.687500)
Diferencia |Bernstein - De Casteljau| = 0.000E+000

==================================================
t = 1.00   (1 - t = 0.00)
==================================================
[Bernstein]  B(t) = suma C(n,i) * (1-t)^(n-i) * t^i * P_i
  i=0: B0,3 = 1 * 0.00^3 * 1.00^0 = 0.000000  ->  B * P0 = (0.000000, 0.000000)
  i=1: B1,3 = 3 * 0.00^2 * 1.00^1 = 0.000000  ->  B * P1 = (0.000000, 0.000000)
  i=2: B2,3 = 3 * 0.00^1 * 1.00^2 = 0.000000  ->  B * P2 = (0.000000, 0.000000)
  i=3: B3,3 = 1 * 0.00^0 * 1.00^3 = 1.000000  ->  B * P3 = (5.000000, 0.000000)
  Suma: B(1.00) = (5.000000, 0.000000)
[De Casteljau]  P_i^(r) = (1-t) * P_i^(r-1) + t * P_(i+1)^(r-1)
  r=0: P0^(0)=(0.000000, 0.000000)  P1^(0)=(1.000000, 3.000000)  P2^(0)=(4.000000, 3.000000)  P3^(0)=(5.000000, 0.000000)  
  r=1: P0^(1)=(1.000000, 3.000000)  P1^(1)=(4.000000, 3.000000)  P2^(1)=(5.000000, 0.000000)  
  r=2: P0^(2)=(4.000000, 3.000000)  P1^(2)=(5.000000, 0.000000)  
  r=3: P0^(3)=(5.000000, 0.000000)  
  Resultado: C(1.00) = (5.000000, 0.000000)
Diferencia |Bernstein - De Casteljau| = 0.000E+000

RESUMEN
  t     | Bernstein                | De Casteljau
  0.00  | (0.000000, 0.000000)     | (0.000000, 0.000000)
  0.25  | (1.062500, 1.687500)     | (1.062500, 1.687500)
  0.50  | (2.500000, 2.250000)     | (2.500000, 2.250000)
  0.75  | (3.937500, 1.687500)     | (3.937500, 1.687500)
  1.00  | (5.000000, 0.000000)     | (5.000000, 0.000000)

Verificacion con Evaluar(lstPuntos, 5) de cada clase:
  Diferencia maxima       = 0.000E+000
  Diferencia promedio     = 0.000E+000
  Error cuadratico medio  = 0.000E+000
\end{lstlisting}

% RESUMEN: Esta seccion (08_anexos.tex) contiene el codigo completo del Modelo
% (lst:anx_bernstein, lst:anx_casteljau, lst:anx_comparador, lst:anx_curva) y la salida
% completa de la validacion manual (lst:anx_salida, subsec:anx_salida). El codigo se copio
% de los archivos reales del repositorio; si el Modelo cambia, se debe actualizar aqui.
% Andres y Joseph pueden agregar el codigo de sus clases en el lugar indicado.
